\guidechapter{Changing the grid without sorting everything}{Read before Gaussian resampling.}

The natural convolution layout and the efficient synthetic-transform layout
do not have the same axis lengths. The multiplication reduction uses pairwise
coprime source axes; synthetic transforms prefer powers of two. Gaussian
resampling connects them through carefully controlled linear maps. This is
an analytic bridge, not ordinary rounding of an index to its nearest neighbor.

In one dimension the source gives an identity of the form
\[
 P_sF_s=2^{2\alpha^2}B_0P_tF_tA.
\]
Here $A$ maps from the source grid of length $s$ to a longer grid of length
$t$; $B_0$ maps back. The $P$ operators are coordinate permutations.
The maps are implemented approximately at a stated precision, with explicit
error and time bounds. The large known scalar is accounted for in recovery.
The inequalities relating $s,t,\alpha,p$ are assumptions of this interface,
not decoration that can be dropped when changing parameters.

\paragraph{Why leaving a permutation in a formula can save time.}
Let $P$ reorder coordinates. If two arrays use the same reordering, then
\[
 (PA)\odot(PB)=P(A\odot B).
\]
Corresponding frequencies are still paired correctly. A consumer that knows
the order can use them directly. The simple identity below explains the
opportunity; the source proves the more specific identities needed to make
all its source and target permutations compatible with the opposite transform.

\explfigure{08-permutations}{Consistency can remove a needless reorder.}{
The arrows show the same permutation on both factors and their pointwise
product. The diagram alone does not prove that an arbitrary inverse-transform
implementation accepts that order. Its interface must establish this.}

\paragraph{Where the chirp identity fits.}
Bluestein's identity rewrites the bilinear exponent using
$2jk=j^2+k^2-(j-k)^2$. Quadratic phase factors on the inputs and outputs
turn a transform into a convolution with a known chirp. The original proof
keeps the relevant permutation inside that calculation. A further phase
twist reconciles cyclic convolution with the negacyclic polynomial relation
$y^r=-1$. Each twist and its removal has both an algebraic and a precision cost.

\paragraph{The error budget has two different kinds of growth.}
For a composition of exact contractions, existing error is not amplified,
so local approximation errors can be added. However, undoing a known large
normalization multiplies absolute error. It is therefore wrong to conclude
``all transforms are contractive, so the final answer is accurate'' without
tracking the rescaling. The source chooses sufficient working precision to
survive both accumulated local error and the final amplification.

\takeaway{A layout is part of an algorithm's input and output contract.
Preserving a useful order is valuable only when every later step accepts it.}
\sourceanchor{07-resampling.tex, lem:no-sort-resampling; 08-assembly.tex}
